\chapter{Data-model reference}
\label{app:datamodel}

Members of the classes in \repo{data-model}, in namespace \texttt{SHiP}. The
first column gives the name from version 0.6, used throughout this booklet;
the second gives the 0.5 name, which you will meet in older code and files.
Not every member changed.
The compatibility limitations in \cref{ch:datamodel} still apply. Positions are in mm, momenta in GeV/$c$, energies in GeV, times
in ns. Arrays of three are $(x, y, z)$.

\newcommand{\dmrow}[4]{\texttt{#2} & \texttt{#1} & #3 & #4 \\}
\newenvironment{dmtable}[1]
  {\par\medskip\setlength{\parfillskip}{0pt plus 1fil}\noindent\begin{minipage}{\linewidth}
   \subsection*{#1}\small\setlength{\tabcolsep}{4pt}
   \setlength{\parfillskip}{0pt plus 1fil}
   \noindent\tabular{@{}P{0.22\linewidth} P{0.22\linewidth} P{0.17\linewidth} P{\dimexpr0.39\linewidth-24pt\relax}@{}}
   \toprule from 0.6 & 0.5 & type & meaning \\ \midrule}
  {\bottomrule\endtabular\par\end{minipage}\par\medskip}

\begin{dmtable}{EventHeader}
\dmrow{weight}{weight}{double}{event weight; 1 unless the generator weights events}
\dmrow{original\_event\_id}{original\_event\_id}{int64}{the event this one came from: for a replay, the event number of the earlier stage; $-1$ if unknown}
\end{dmtable}

\begin{dmtable}{MCParticle}
\dmrow{pdgCode}{pdg\_code}{int32}{PDG particle code}
\dmrow{vertex}{vertex}{double[3]}{production vertex}
\dmrow{momentum}{momentum}{double[3]}{momentum}
\dmrow{energy}{energy}{double}{\textbf{total} energy}
\dmrow{time}{time}{double}{production time}
\dmrow{motherId}{mother\_id}{int32}{index of the first mother in the same vector; $-1$ for none}
\dmrow{mothers}{mothers}{vector<int32>}{indices of all mothers (from 0.5)}
\dmrow{status}{status}{int32}{HepMC status; 1 = tracked by Geant4. Plane particles: 1 for the tracks at the plane, others for their ancestors}
\end{dmtable}

\begin{dmtable}{SimParticle}
\dmrow{trackId}{track\_id}{int32}{Geant4 track ID, unique within an event}
\dmrow{parentId}{parent\_id}{int32}{track ID of the creator; 0 for particles from the source}
\dmrow{pdgCode}{pdg\_code}{int32}{PDG particle code}
\dmrow{vertex}{vertex}{double[3]}{production vertex}
\dmrow{endpoint}{endpoint}{double[3]}{where tracking stopped}
\dmrow{momentum}{momentum}{double[3]}{initial momentum}
\dmrow{energy}{energy}{double}{initial \textbf{kinetic} energy}
\dmrow{time}{time}{double}{production time}
\dmrow{creatorProcess}{creator\_process}{int32}{Geant4 process sub-type that created the track}
\end{dmtable}

\begin{dmtable}{SimHit}
\dmrow{detectorId}{detector\_id}{int32}{0 UBT, 1 SBT, 2 straws, 3 calorimeter, 4 timing; 100 for an \pkg{aegir} scoring plane}
\dmrow{geometryNodeId}{geometry\_node\_id}{int32}{the scoring plane's id for plane hits; 0 for detector hits}
\dmrow{trackId}{track\_id}{int32}{Geant4 track that made the hit}
\dmrow{pdgCode}{pdg\_code}{int32}{PDG code of that track}
\dmrow{position}{position}{double[3]}{hit position}
\dmrow{momentum}{momentum}{double[3]}{momentum at the hit}
\dmrow{energyDeposit}{energy\_deposit}{double}{energy deposited in the step; 0 for plane and crossing-mode hits}
\dmrow{time}{time}{double}{global time}
\dmrow{pathLength}{path\_length}{double}{step length}
\end{dmtable}

\begin{dmtable}{RecHit, and the detector hits wrapping it}
\dmrow{detectorId}{detector\_id}{int32}{as in \texttt{SimHit}}
\dmrow{trackId}{track\_id}{int32}{truth link, as in \texttt{SimHit}}
\dmrow{pdgCode}{pdg\_code}{int32}{truth link, as in \texttt{SimHit}}
\dmrow{position}{position}{double[3]}{measured position}
\dmrow{momentum}{momentum}{double[3]}{as in \texttt{SimHit}}
\dmrow{energyDeposit}{energy\_deposit}{double}{as in \texttt{SimHit}}
\dmrow{time}{time}{double}{measured time}
\dmrow{pathLength}{path\_length}{double}{as in \texttt{SimHit}}
\multicolumn{4}{@{}p{\linewidth}@{}}{\texttt{UBTHit}, \texttt{SBTHit}, \texttt{StrawTubesHit}, \texttt{CaloHit},
  \texttt{TimeDetHit}: one member \texttt{rec\_hit} (0.5: \texttt{recHit}).} \\
\end{dmtable}

\begin{dmtable}{RecParticle (reconstruction)}
\dmrow{trackId}{track\_id}{int32}{}
\dmrow{parentId}{parent\_id}{int32}{}
\dmrow{pdgCode}{pdg\_code}{int32}{}
\dmrow{vertex}{vertex}{double[3]}{production vertex}
\dmrow{endpoint}{endpoint}{double[3]}{end point}
\dmrow{momentum}{momentum}{double[3]}{initial momentum}
\dmrow{hits}{hits}{vector<RecHit>}{hits on the particle}
\dmrow{energy}{energy}{double}{initial kinetic energy}
\dmrow{time}{time}{double}{production time}
\dmrow{creatorProcess}{creator\_process}{int32}{}
\dmrow{ipPV}{ip\_pv}{double}{impact parameter w.r.t.\ the primary vertex (now: the origin)}
\end{dmtable}

\begin{dmtable}{TrackFitResult (reconstruction)}
\dmrow{nMeas}{n\_meas}{int32}{number of measurements}
\dmrow{fitStatus}{fit\_status}{int32}{0 success, 1 failure}
\dmrow{chi2}{chi2}{double}{}
\dmrow{ndf}{ndf}{int32}{}
\dmrow{qoverp}{q\_over\_p}{double}{charge over momentum}
\dmrow{phi}{phi}{double}{}
\dmrow{theta}{theta}{double}{}
\dmrow{time}{time}{double}{}
\dmrow{refLoc}{ref\_loc}{double[3]}{reference location}
\multicolumn{4}{l}{plus input and fitted measurements and residuals in $x$ and $y$ (vectors).} \\
\end{dmtable}
